\section{From local algorithms to a global complexity argument}
\label{sec:global}

\subsection{An explicit ledger of the remaining obligations}

A local algorithm can be fixed-parameter tractable while the enclosing
recognizer remains exponential.  There are four independent quantities to
control: the cost of producing a legal next surface or move, the number of
states visited, the encoded size of each state, and the cost of resolving
the terminal topology.  A small matrix at one successful cut gives no bound
on the other three quantities.

The reviewed hierarchy preprint supplies a polynomial compressed-cutting
primitive for handle structures of uniform type, polynomial in handle count
and logarithmic extended weight.  It returns the complement of the parallelity
bundle and the bundle topology and attachment data (Theorem 14.4), while
Section 9 separates a step-count analysis from
a full runtime estimate \cite{Lackenby2026}.  This is the appropriate
interface for the present work: geometry must arrive in a compressed,
source-authenticated representation, and its attachment information must
survive every simplification.

\begin{center}
\begin{tabular}{P{0.25\textwidth}P{0.31\textwidth}P{0.33\textwidth}}
\toprule
Obligation & Result available here & Still required globally\\
\midrule
Raw local completion & Exact $J$-state reduction; deterministic width
algorithm & Completeness of the chosen geometric move class at every
needed state.\\
Representation size & No expansion in multiplicity in these finite
operators & Bounds on handles, charts, runs, and integer bit lengths
after all cuts and transports.\\
Hierarchy control & Explicit local cost parameters & A bound on total
states, restarts, and surface-production work.\\
Boundary and bundle data & Existing native verifier; five new accepted
packets & Complete primitive-attachment, puncture, and terminal handling
where current certificates refuse.\\
Randomness & One-sided local witness misses, exact positive replay & A
global error schedule, or a deterministic structural regime.\\
\bottomrule
\end{tabular}
\end{center}

\subsection{A conditional composition bound}

The following accounting statement makes the target testable.  It is a
conditional composition result, not a theorem that the hypotheses hold
for all knot diagrams.

\begin{proposition}[A sufficient deterministic route to quasi-polynomial time]
\label{prop:global-conditional}
Suppose a correct and complete recognition algorithm on $n$-crossing input
diagrams has the following independently established properties:
\begin{enumerate}
\item it visits at most $Q(n)$ states and makes at most polynomially many
  local completion calls per state;
\item each call is equivalent to the finite completion problem of this
  article, with explicit encoded size $N\leq n^b$ for a fixed $b$;
\item all remaining production, transport, and verification work per state
  is bounded by $n^b$ after increasing $b$ if necessary;
\item the implication width of every local call is at most $w(n)$.
\end{enumerate}
If $Q(n)\leq\exp(O((\log n)^a))$ and
$w(n)=O((\log n)^c)$ for fixed $a,c\geq0$, then replacing its finite
completion calls by the deterministic width algorithm gives total time
\[
       \exp\!\left(O\left((\log n)^{\max\{a,c+1,1\}}\right)\right).
\]
In particular it is quasi-polynomial.  The correctness and completeness
assumption includes all geometric and terminal decisions, not only the
existence of the local certificates considered here.
\end{proposition}
\begin{proof}
For $w>0$, the ideal bound gives $J\leq(d+1)^w\leq(N+1)^w$.
The deterministic width algorithm therefore costs $N^{O(w+1)}$ per call.
For $w=0$ it is polynomial as well.  Substituting $N\leq n^b$ gives
\(\exp(O((w(n)+1)\log n))\).  Multiplying by the permitted number of
states and calls and adding the other work yields the displayed bound.
Because each local algorithm solves exactly the specified finite problem,
this replacement does not alter an independently proved global correctness
argument.
\end{proof}

The randomized alternative replaces the local factor by
$2^d\poly(N)$, or more precisely $J\poly(N)$.  With a known upper bound
on the number of relevant calls, distribute a desired total miss allowance
$\varepsilon$ among them and repeat independent identity tests accordingly.
A union bound controls the probability of any required local witness being
missed, provided all other completeness obligations are met.  A zero test
still cannot be treated as a deterministic negative certificate.

For example, polynomial state size and $d=O(\log n)$ give a polynomial
randomized local cost.  Polynomial state size and $w=O(\log n)$ give a
deterministic local cost $\exp(O((\log n)^2))$.  If a hierarchy visits
$\exp(O((\log n)^3))$ states, the total bound remains quasi-polynomial
but has the larger exponent.  This explains why an article must specify
whether it controls local work, the state count, or both.

The bound $k\leq2T$ for a raw core with $T$ tetrahedra is insufficient on
its own: $8^{2T}$ is exponential in $T$.  Nor does observing
$\kappa\leq1$ on 106 inputs establish such a bound on an infinite class.
The useful advance is that there are now precise quantities---$\kappa$,
$d$, $J$, $w$, and $q$---for a global geometric theorem to control.

\section{Proposed research program}
\label{sec:research}

The next projects should have a definite theorem, failure example, or
replayable artifact as their output.  The order below emphasizes the two
current positive gaps in annular attachment, then the structural questions
most directly connected to asymptotic complexity.

\subsection{Geometric continuation and choice of lower forest}

\paragraph{1. Complete the two newly exposed Fibonacci attachment cases.}
The index-eight and index-ten Fibonacci cuts now admit raw core
cancellation under the optimized lower forest, yet their annular interface
stage refuses.  Export the exact remaining annular curves, the lower-forest
basis, and all source endpoint maps.  Determine whether a verified change
of meridian system makes the attachment primitive in the required geometric
sense.  A useful result is a family theorem with a compressed construction,
not just two extra successes.  The deliverable should be a source-bound
move sequence followed by an unchanged terminal check, together with a
bound in layer count and coordinate bit length.

\paragraph{2. Optimize lower forests for completion, not only for $k$.}
The graphic-nullity formula solves the parameter objective exactly, but
minimizing that objective need not preserve a successful attaching pattern.
Study the joint choice of lower and top forests.  First search for an
explicit native example where every minimum-$k$ lower forest fails but
another lower forest succeeds.  If such examples exist, identify a secondary
objective or a bounded family of exchanges that preserves completeness.
If they do not occur in a well-defined native subclass, prove the exchange
property.  The reverse-order rescue of the Pachner case already shows why
a single numerical score should not silently replace a portfolio.

\paragraph{3. Authenticate geometric tightening and handle slides.}
Actual unsigned occurrences can exceed the algebraic obstruction after
curves are tightened or the meridian system changes.  Signed exponent sums,
gcd tests, or free-group simplification alone do not provide the geometric
disc required by the current certificate.  Develop a compiler for an
embedded boundary-curve isotopy or a geometrically realized handle slide,
with certified endpoint transport and intersection counts.  Existing
polynomial compressed-curve algorithms supply relevant primitives
\cite{LackenbyCurves2026}; the open work is their integration with this
marked source representation and simultaneous annular attachments.
The first milestone is a checker for one move, including its inverse and
the preservation of boundary pattern and orientability.

\paragraph{4. Specify punctured-handlebody terminals separately.}
Six original refusals contain non-handlebody pieces according to the
historical oracle.  Some applications may instead want a handlebody after
certified sphere capping.  Define that terminal class explicitly, including
which boundary spheres are capped and how the original component is
recovered.  A new puncture certificate should not be accepted as an ordinary
handlebody certificate.  This project can turn otherwise misleading
``misses'' into a useful terminal specification while preserving the
logical meaning of the existing genus census.

\subsection{Sharper structural parameters and exact algorithms}

\paragraph{5. Exploit native cycle compatibility of triple supports.}
The current theorem permits arbitrary exceptional supports.  Source
boundary compatibility imposes additional mod-two cycle conditions in the
dual graph, as audited in the preceding normalization report
\cite{IncomingForests2026}.  A three-edge even subgraph has very restricted
forms, involving triangles, parallel edges, and loops.  Determine whether
these constraints give a stronger exchange theorem, a smaller effective
parameter, or a polynomial class beyond support two.  The target is an
exact reduction respecting actual unsigned entries; parity can constrain
the search but cannot substitute for a unit crossing.  Any hardness
construction used to delimit this problem must satisfy the native source
constraints before being called a topological hardness result.

\paragraph{6. Saturate unary consequences before building the poset.}
Forward chaining from one row may imply another through several nonlinear
rules.  Such a unary logical consequence is safe to add as an ordinary
implication if its derivation is retained.  Investigate whether repeatedly
adding these consequences substantially collapses the special-SCC poset,
and prove the exact correspondence with the original closure semantics.
The obvious polynomial preprocessing tries closure from each row.  The
interesting questions are whether a smaller set of representatives suffices,
how much $J$ or $w$ can decrease, and whether native zero patterns impose
a uniform bound on the saturated width.  The positive checker can continue
to replay the final tree and pivots in the original matrix.

\paragraph{7. Determine the best exact route when $q$, $w$, and $d$ differ.}
The present alternatives have costs $J\poly(N)$ with one-sided randomness,
$J(m+1)^w\poly(N)$ deterministically, and $2^qJ\poly(N)$ by exact counting.
Implement the Kirchhoff alternative with fraction-free or certified modular
determinants, then compare these regimes using source-derived parameters.
Seek a deterministic $2^{O(d)}\poly(N)$ algorithm when $q$ is large.
The signed common-basis polynomial is the concrete difficulty; a positivity
argument for a different polynomial must prove that it has exactly the
required support.  General nonnegative-polynomial detection results are
relevant context, but do not directly apply to the mixed-sign common-basis
enumerator \cite{BrandKorchemnaSkotnica2023}.

\paragraph{8. Replace large explicit ideal tables when the frontier is small.}
A poset can have many ideals even when a particular elimination ordering
has a small active frontier.  Study a decision diagram or a dynamic program
that compiles the activation function and ideal constraints together.
Specify exactly what boundary information must be retained so that a tree
outside the processed region is represented faithfully.  It is not enough
to compress the Boolean function if the graphic common-basis constraints
lose their dependence across the boundary.  A useful theorem would express
time in an explicitly defined width of the combined dependency and graphic
structure, with an independent witness extractor.

\subsection{Verification, representation size, and the hierarchy}

\paragraph{9. Produce portable negative finite certificates.}
The deterministic search can exhaust its finite space, but most of its
negative outputs currently have no exported proof object.  Develop a
certificate that an independent checker can validate without trusting the
search program.  Possible starting points are per-state graphic rank
obstructions, a shared acyclic proof of all failed representative choices,
or exact Kirchhoff filtered counts in the small-$q$ regime.  The certificate
must remain bound to the original graph and matrix, and its verification
cost and size should be included in the parameter bound.  Randomized zero
evaluations are insufficient.

\paragraph{10. Improve the remaining compressed transport costs.}
The prefix theorem removes repeated gap sums, but the measured runtime
still includes weight-run overlays, translations, reflection folds,
selector restoration, and repeated trace traversals.  Profile these after
the current patch.  Investigate sorted batched endpoint transport and
interval-run composition, with explicit bit bounds and retained/deleted
endpoint conventions.  A proposed improvement should preserve independently
authenticated gap lists and demonstrate a gain on complete native replay.
The corpus's $2^{16384}$ coordinate scale is a useful check that a convenient
array representation has not introduced point expansion.

\paragraph{11. Build native stress families for the actual new parameters.}
The present raw corpus has dual rank at most one and very small ideal
counts.  Create triangulated source families with increasing dual rank,
controlled exceptional-column cycles, nontrivial implication chains, and
large binary coordinates.  Report $T$, row count, dual rank, $k$, $\kappa$,
$d$, $J$, $w$, $q$, certificate size, and replay time.  Keep abstract
graph/matrix generators separately labelled: they are excellent logical
tests but do not establish realizability by a triangulated source.
Prioritize families that distinguish an exponential number of activation
masks from a polynomial number of ideals.

\paragraph{12. Prove a hierarchy-wide complexity budget.}
Track the number of active nonparallel handles, boundary-pattern records,
interval runs, restart events, and coordinate bit lengths throughout a
complete hierarchy implementation.  Seek an amortized invariant bounding
both the number of visited states and the maximum local ideal complexity.
A decrease in handle count alone is insufficient if a restart duplicates
many states or an attachment representation grows too quickly.  The desired
deliverable is a theorem supplying the hypotheses of
\cref{prop:global-conditional}, or a counterexample showing which proposed
potential fails.  This is the central remaining bridge from the proved
local operators to Lackenby's announced asymptotic target.

\subsection{Recommended next integration checkpoint}

The behavior-preserving speed patch is ready for ordinary code review and
the repository's usual integration gates.  The lower-forest portfolio is
separate and can be reviewed with its five full packets and shared resource
limits.  The exceptional-row solvers should enter as research operators
with measured dispatch parameters, explicit state/work limits, and the
unchanged finite verifier.  Their source adapter must carry all lower-forest
and attachment provenance before they are used as native geometric claims.

For the next research iteration, the highest-value concrete target is the
source-authenticated annular continuation of the two larger Fibonacci
examples.  In parallel, the structural priority is to test and then prove
a bound on implication width after geometric simplification.  Those tasks
would connect the native coverage improvement and the new deterministic
complexity theorem at an interface where both can be checked precisely.
